\section{Introduction}
\label{section:intro}
%Bakgrund och context
%Problemet, vad vill vi åtgärda
Sign language is a visible language expressed through hand gestures, eye gaze, facial expression and other small visible cues \cite{Sign_lang}. Today, sign language is spoken by more than 72 million people worldwide. However, the difficulty for those who use and rely on sign language to communicate with those who do not know it remains challenging, as is common in situations where individuals attempt to communicate across language barriers. Text-to-speech and transcription tools can partially bridge this gap but being unable to use their usual method of communication can be taxing. There are several commercially available experimental projects that demonstrate functional sign-to-speech gloves that aim to alleviate this problem.

%Relaterat arbete
\subsection{Related Works}
One notable example is BrightSign, which emphasizes a wide library of supported languages, customizable voice options, and a companion app that enhances the glove's usability at a steep price of 27031 sek at the time of writing\cite{bright}.

Another significant development was presented by UCLA researchers in 2019, who created a prototype capable of recognizing American Sign Language (ASL) gestures, aiming for real-time translation into text and speech \cite{UCLA}.

%Vart ligger vårt fokus
Although these projects demonstrate the feasibility of translating sign language through wearable technology, they have a high costs, expensive components or require specialized hardware and software ecosystems.